% !PRINTED 267
\weberpage{267}
Wir ordnen die Determinante so an:
\[
\begin{vmatrix}
-2a, & -3b, & -4c \\
-3b, & a^2 - 4c, & \tfrac{1}{2} a b \\
-4c, & \tfrac{1}{2} a b, & -\tfrac{1}{4} (8ac - 3b^2)
\end{vmatrix}
\]
und erhalten die Kette der Haupt-Unterdeterminanten
\weq{3}{D, \; -2a^3 + 8ac - 9b^2, \; -2a, \; 1,}
wenn $D$ die Discriminante bedeutet, die nach \S.\,64, (11) bestimmt ist durch
\weq{4}{27 D = 4(a^2 + 12c)^3 - (2a^3 - 72ac + 27b^2)^2.}
Das Kennzeichen dafür, dass alle vier Wurzeln reell sind, ist hiernach
\weq{5}{D > 0, \quad -2a^3 + 8ac - 9b^2 > 0, \quad -2a > 0.}
Dies Kennzeichen ist scheinbar verschieden und jedenfalls weniger einfach, als das in \S.\,78 aufgestellte
\weq{6}{D > 0, \quad a < 0, \quad a^2 - 4c > 0.}

Wir wollen aber nun noch nachweisen, dass beides genau dasselbe besagt.

Zunächst ist sofort zu übersehen, dass (6) erfüllt sein muss, wenn (5) erfüllt ist, denn
\[
- 2a(a^2 - 4c) - 9b^2
\]
kann nicht positiv sein, wenn $a$ und $a^2 - 4c$ negativ sind; um aber umgekehrt einzusehen, dass aus (6) die Bedingungen (5) folgen, müssen wir den Werth von $D$ in Betracht ziehen.

Wir setzen zur Abkürzung
\weq{7}{\alpha = -2a, \quad \beta = -6a^3 + 24ac - 27b^2, \quad \gamma = 3a^2 - 12c,}
so dass die Bedingungen (5)
\[
\alpha > 0, \quad \beta > 0, \quad D > 0,
\]
die Bedingungen (6)
\[
\alpha > 0, \quad \gamma > 0, \quad D > 0
\]
lauten; dann ist nachzuweisen, dass, wenn $\alpha$, $\gamma$ und $D$ positiv sind, auch $\beta$ positiv sein muss. Hierzu drücken wir $D$ in (4) nach (7) durch $\alpha, \beta, \gamma$ aus und erhalten
\[
27 D = 4(\alpha^2 - \gamma)^3 - [2\alpha(\alpha^2 - \gamma) - \beta]^2.
\]
Wenn $D$ positiv ist, so muss hiernach gewiss $\alpha^2 - \gamma$ positiv sein. Setzen wir also
%CONT>>
